\documentclass[11pt,letterpaper]{article}

\usepackage[margin=1in]{geometry}
\usepackage{booktabs}
\usepackage{xcolor}
\usepackage{tcolorbox}
\usepackage{fancyhdr}
\usepackage{enumitem}
\usepackage{titlesec}
\usepackage{array}
\usepackage{tabularx}
\usepackage{hyperref}
\usepackage{multicol}

\definecolor{projblue}{RGB}{0,70,127}
\definecolor{projgreen}{RGB}{0,100,50}
\definecolor{warnred}{RGB}{180,60,0}
\definecolor{lightgray}{RGB}{248,248,248}
\definecolor{donegray}{RGB}{200,200,200}

\titleformat{\section}{\large\bfseries\color{projblue}}{\thesection.}{0.5em}{}[\vspace{-4pt}\rule{\linewidth}{0.5pt}]
\titleformat{\subsection}{\normalsize\bfseries\color{projgreen}}{\thesubsection}{0.5em}{}
\titleformat{\subsubsection}{\normalsize\bfseries}{\thesubsubsection}{0.5em}{}

\pagestyle{fancy}
\fancyhf{}
\lhead{\textcolor{projblue}{\textbf{Research Handover Document — Binghamton University}}}
\rhead{\textcolor{gray}{April 2026}}
\cfoot{\thepage}
\renewcommand{\headrulewidth}{0.4pt}

\setlist[itemize]{itemsep=3pt, topsep=4pt}
\setlist[enumerate]{itemsep=3pt, topsep=4pt}

\begin{document}

% ── Cover ────────────────────────────────────────────────────────────────────
\begin{center}
{\LARGE\bfseries\color{projblue} Research Project Handover Document}\\[6pt]
{\large Binghamton University · Department of Systems Science \& Industrial Engineering}\\[4pt]
{\normalsize Prepared by: Sriirohit Prakash \quad|\quad \texttt{sprakash1@binghamton.edu}}\\[4pt]
{\normalsize April 2026}\\[10pt]
\rule{\linewidth}{1pt}
\end{center}

\vspace{0.3cm}
\begin{tcolorbox}[colback=projblue!6, colframe=projblue,
  title=\textbf{Overview — Two Active Research Projects}]
This document provides a complete handover for two parallel research projects under
Professor Ertem's supervision. Each project has its own dataset, codebase, and research
objective. Both are in active development and require continued work.

\medskip
\begin{tabular}{llll}
\toprule
\textbf{Project} & \textbf{Domain} & \textbf{Data} & \textbf{Status} \\
\midrule
Project 1 & Colon Cancer Metastasis Prediction & SEER 21 Registry & Model built; validation pending \\
Project 2 & ADRD \& Fall Hospitalizations       & HCUP NIS 2016–2022 & Cohorts built; Table 1 in progress \\
\bottomrule
\end{tabular}
\end{tcolorbox}

\newpage

% ══════════════════════════════════════════════════════════════════════════════
\section{Project 1 — Colon Cancer Metastasis Prediction (SEER ML)}
% ══════════════════════════════════════════════════════════════════════════════

\subsection{Vision and Research Question}

\begin{tcolorbox}[colback=lightgray, colframe=projblue!50]
\textbf{Core question:} Can we predict the probability of distant metastasis in colon
cancer \emph{before} comprehensive imaging, using only information available at the
moment of initial diagnosis?

\medskip
\textbf{Clinical goal:} Provide an early risk stratification tool that helps oncologists
estimate cure rate and longevity, enabling faster treatment planning without waiting for
full staging workup. This is especially valuable in resource-limited settings.
\end{tcolorbox}

\subsection{Dataset Required}

\begin{tcolorbox}[colback=warnred!8, colframe=warnred,
  title=\textbf{Dataset: SEER 21 Registry (National Cancer Institute)}]
\begin{itemize}
  \item \textbf{Source:} Surveillance, Epidemiology, and End Results (SEER) Program —
        21 registry version, available at \url{https://seer.cancer.gov}
  \item \textbf{Access:} Requires SEER Research Data Agreement. Apply at SEER website;
        approval typically takes 2–4 weeks.
  \item \textbf{Files needed:}
        \texttt{export.txt} (main data) + \texttt{export.dic} (data dictionary)
        — exported from SEER*Stat software with colon cancer filters
  \item \textbf{Filter applied in SEER*Stat before export:}
        Site = Colon (C18.1–C18.9); Years = 2000–2022; Histology = malignant
  \item \textbf{Size:} $\sim$809,625 patients after our exclusions
  \item \textbf{Key variables used:} Age at diagnosis, Sex, Race, Primary site (ICD-O-3),
        Histology (ICD-O-3 morphology code), Tumor grade, Year of diagnosis,
        SEER Summary Stage (metastasis label — our target variable)
  \item \textbf{NOT available in SEER (limits model):} T/N staging details, CEA levels,
        tumor size, treatment type — these will come from Guthrie dataset (Phase 2)
\end{itemize}
\end{tcolorbox}

\subsection{What Has Been Done}

\subsubsection{Data Preprocessing (\texttt{seer\_model.py})}
\begin{itemize}
  \item Loaded raw SEER fixed-width export file (\texttt{export.txt})
  \item Applied inclusions: Grades I, II, III only; colon sites C18.1–C18.9;
        years 2000–2022
  \item Applied exclusions:
    \begin{itemize}
      \item Grade IV (undifferentiated/anaplastic, $\sim$9,999 pts) —
            biologically distinct, distorts predictions
      \item Immunophenotype labels (B-cell, T-cell, NK cell, Null cell,
            $\sim$3,670 pts) — immune lineage labels, not valid cancer grades;
            appear in SEER due to multi-cancer encoding
    \end{itemize}
  \item Created binary label: \texttt{Metastasis} = 1 if SEER Summary Stage = ``Distant'',
        else 0
  \item One-hot encoded categorical variables with \texttt{drop\_first=True}
  \item 70/30 stratified train/test split
\end{itemize}

\subsubsection{Model Training (\texttt{seer\_model.py})}
\begin{itemize}
  \item Baseline: Logistic Regression with StandardScaler
  \item Final model: XGBoost with \texttt{CalibratedClassifierCV}
        (Platt sigmoid calibration) — calibration converts raw XGBoost scores
        into reliable probabilities
  \item \textbf{Important:} StandardScaler is only applied to Logistic Regression.
        XGBoost is trained on raw (unscaled) data. Do NOT apply scaler to XGBoost
        inference — this was a past bug that caused all predictions to return $\sim$1.2\%.
  \item Saved artifacts: \texttt{xgb\_model.pkl}, \texttt{scaler.pkl},
        \texttt{X\_train.columns} (column order must be preserved for inference)
\end{itemize}

\subsubsection{Model Performance}
\begin{center}
\begin{tabular}{lcccc}
\toprule
\textbf{Metric} & \multicolumn{2}{c}{\textbf{Logistic Regression}} &
\multicolumn{2}{c}{\textbf{XGBoost (final)}} \\
\cmidrule(lr){2-3}\cmidrule(lr){4-5}
 & Non-Met & Metastatic & Non-Met & Metastatic \\
\midrule
AUC-ROC   & \multicolumn{2}{c}{0.673} & \multicolumn{2}{c}{\textbf{0.785}} \\
Precision & 0.83 & 0.33 & 0.85 & 0.61 \\
Recall    & 1.00 & 0.00 & \textbf{0.98} & 0.19 \\
F1-Score  & 0.91 & 0.00 & 0.91 & 0.29 \\
Accuracy  & \multicolumn{2}{c}{0.83} & \multicolumn{2}{c}{\textbf{0.84}} \\
\bottomrule
\end{tabular}
\end{center}

\medskip
\noindent\textbf{Primary focus — Non-Metastatic recall (98\%):} The model correctly
identifies 98\% of non-metastatic patients. Metastatic recall is low (19\%) because
the strongest metastasis predictors (T/N stage, CEA, tumor size) are not in SEER.

\subsubsection{Inference Pipeline (\texttt{seer\_predict.py})}
\begin{itemize}
  \item Load new patient data as a DataFrame
  \item One-hot encode with \texttt{drop\_first=False} (different from training!)
  \item Align columns: \texttt{reindex(columns=X\_train.columns, fill\_value=0)}
        — this is critical; misalignment was the root cause of the 1.2\% bug
  \item Apply XGBoost model directly (no scaler)
  \item Output: \texttt{Metastasis\_Probability} and \texttt{NonMetastasis\_Probability}
\end{itemize}

\subsubsection{Streamlit Web App (\texttt{app.py})}
\begin{itemize}
  \item Interactive UI for single-patient prediction
  \item Inputs: Age group, Sex, Race, Primary site, Histology (ICD-O-3), Grade, Year
  \item Outputs: Metastasis probability (\%), Non-metastasis probability (\%),
        Risk level (Low/Moderate/High)
  \item Grade dropdown: Grade I, II, III only (Grade IV removed)
  \item Run with: \texttt{streamlit run app.py}
\end{itemize}

\subsubsection{Report}
\begin{itemize}
  \item \texttt{SEER\_ML\_Report\_v2.tex} — 3-page LaTeX report with objective,
        inclusions/exclusions, model accuracy table, 10 test cases, and screenshots
  \item Needs: \texttt{out1.png}–\texttt{out10.png} (Streamlit screenshots) uploaded
        to Overleaf to render Section 5 (screenshots)
\end{itemize}

\subsection{What Still Needs to Be Done}

\begin{tcolorbox}[colback=warnred!8, colframe=warnred,
  title=\textbf{Phase 2 — Guthrie Dataset External Validation (PENDING)}]
\begin{enumerate}
  \item \textbf{Obtain Guthrie dataset} — clinical records from Guthrie Health System.
        Contains variables absent from SEER: T/N stage, CEA levels, tumor size,
        treatment type. This is the external validation dataset.
  \item \textbf{External validation} — run \texttt{seer\_predict.py} on Guthrie patients
        and compare predicted probabilities to actual outcomes.
  \item \textbf{Expected improvement} — metastatic recall projected to improve from
        19\% to $>$50\% once T/N staging and CEA are added as features.
  \item \textbf{Retrain if needed} — if Guthrie data reveals new variables, retrain
        XGBoost with expanded feature set.
\end{enumerate}
\end{tcolorbox}

\subsection{Key Files — Project 1}

\begin{center}
\begin{tabular}{lp{9cm}}
\toprule
\textbf{File} & \textbf{Purpose} \\
\midrule
\texttt{seer\_model.py}       & Data loading, preprocessing, model training, saves .pkl files \\
\texttt{seer\_predict.py}     & Batch inference on new patients \\
\texttt{app.py}               & Streamlit web UI for single-patient prediction \\
\texttt{xgb\_model.pkl}       & Trained XGBoost model (gitignored — do not commit) \\
\texttt{scaler.pkl}           & StandardScaler for Logistic Regression only (gitignored) \\
\texttt{SEER\_ML\_Report\_v2.tex} & Final report LaTeX source \\
\texttt{.gitignore}           & Excludes large data files and model artifacts from GitHub \\
\bottomrule
\end{tabular}
\end{center}

\subsection{How to Run — Project 1}

\begin{enumerate}
  \item Place \texttt{export.txt} and \texttt{export.dic} in the project folder
  \item \texttt{python seer\_model.py} — trains model, saves \texttt{xgb\_model.pkl}
  \item \texttt{streamlit run app.py} — launches web app at \texttt{localhost:8501}
  \item \texttt{python seer\_predict.py} — batch predictions on new patients
\end{enumerate}

\newpage

% ══════════════════════════════════════════════════════════════════════════════
\section{Project 2 — ADRD \& Fall Hospitalizations (HCUP NIS)}
% ══════════════════════════════════════════════════════════════════════════════

\subsection{Vision and Research Question}

\begin{tcolorbox}[colback=lightgray, colframe=projblue!50]
\textbf{Core question:} Among patients with Alzheimer's Disease and Related Dementias
(ADRD) hospitalized in the US, what are the demographic, clinical, and spatial differences
between those admitted \emph{for} a fall, those with a \emph{prior fall history}, and those
with \emph{no fall}?

\medskip
\textbf{Clinical goal:} Understand the burden of fall-related hospitalizations in the
ADRD population at a national level (2016–2022) and identify risk factors that distinguish
the three cohorts. Findings will inform fall prevention strategies for dementia patients.
\end{tcolorbox}

\subsection{Dataset Required}

\begin{tcolorbox}[colback=warnred!8, colframe=warnred,
  title=\textbf{Dataset: HCUP National Inpatient Sample (NIS), 2016–2022}]
\begin{itemize}
  \item \textbf{Source:} Healthcare Cost and Utilization Project (HCUP), AHRQ.
        Available at \url{https://www.hcup-us.ahrq.gov/nisoverview.jsp}
  \item \textbf{Access:} Requires HCUP Data Use Agreement (DUA) — separate agreement
        per year. Annual fee applies. Apply at \url{https://www.hcup-us.ahrq.gov/tech_assist/dua.jsp}
  \item \textbf{Files needed per year:}
    \begin{itemize}
      \item \textbf{Core file} (\texttt{NIS\_YYYY\_Core.ASC}) — patient-level discharge
            records; fixed-width ASCII format; $\sim$7M discharges/year nationally
      \item \textbf{Hospital file} (\texttt{NIS\_YYYY\_Hospital.ASC}) — hospital
            characteristics (region, teaching status, bed size); keyed by \texttt{HOSPID}
            — \textbf{NOT YET MERGED}
    \end{itemize}
  \item \textbf{Column positions:} Vary by year. Defined per-year in
        \texttt{get\_col\_dict()} in \texttt{hcup\_adrd\_falls\_cohorts.py}.
        Always verify against official SAS load files before changing.
  \item \textbf{Critical format change:} NIS redesigned the record layout in 2019.
        Years 2016–2018 use different column positions than 2019–2022.
        Our code handles this correctly.
  \item \textbf{Important limitation:} NIS is \textbf{discharge-level}, not
        patient-level. \texttt{KEY\_NIS} identifies a hospitalization, not a person.
        The same patient can appear multiple times across years with no linkage.
        All cohort counts represent hospitalizations, not unique patients.
\end{itemize}
\end{tcolorbox}

\subsection{Cohort Definitions}

Three mutually exclusive cohorts, applied after exclusions:

\begin{center}
\renewcommand{\arraystretch}{1.4}
\begin{tabularx}{\textwidth}{lXr}
\toprule
\textbf{Cohort} & \textbf{Rule} & \textbf{N (2016–2022)} \\
\midrule
\textbf{ADRD\_FALL\_PRIMARY}
  & ADRD in any DX \textbf{AND} \texttt{I10\_DX1} (principal diagnosis) =
    fall-related injury (S00–S99, T00–T34) or fall mechanism (W00–W19)
  & 214,140 \\
\textbf{ADRD\_PRIOR\_FALL}
  & ADRD in any DX \textbf{AND} Z91.81 (history of falling) in any DX
    \textbf{AND} not in Cohort 1
  & 127,242 \\
\textbf{ADRD\_NO\_FALL}
  & ADRD in any DX; NO fall codes (W00–W19), NO Z91.81 anywhere
  & 2,009,773 \\
\midrule
\textit{UNCLASSIFIED}
  & Fall W-code in secondary DX only; DX1 is not an injury code
    (e.g., admitted for pneumonia with incidental fall code) —
    \textbf{awaiting professor decision on classification}
  & 144,643 \\
\midrule
\textbf{Total} & & \textbf{2,495,798} \\
\bottomrule
\end{tabularx}
\end{center}

\subsection{ADRD ICD-10 Codes Used}

\begin{center}
\begin{tabular}{ll}
\toprule
\textbf{Code} & \textbf{Condition} \\
\midrule
G30.x  & Alzheimer's disease \\
G31.0x & Frontotemporal dementia \\
G31.83 & Dementia with Lewy bodies \\
F01.x  & Vascular dementia \\
F02.x  & Dementia in other diseases \\
F03.x  & Unspecified / general dementia \\
G23.1  & Progressive supranuclear palsy \\
\bottomrule
\end{tabular}
\end{center}

\subsection{Exclusion Criteria}
\begin{itemize}
  \item Age $<$ 50 at admission (\texttt{AGE} variable)
  \item Any malignancy code C00–C96 in any DX field (all 40 diagnosis positions)
\end{itemize}

\subsection{What Has Been Done}

\subsubsection{Cohort Builder (\texttt{hcup\_adrd\_falls\_cohorts.py})}
\begin{itemize}
  \item Reads raw NIS Core ASC files year by year in 100,000-row chunks
  \item Applies ADRD filter, age exclusion, cancer exclusion
  \item Assigns each discharge to one of the three cohorts using vectorized regex
  \item Computes 10 Elixhauser comorbidity flags per patient using ICD-10 regex
  \item Saves results to \texttt{HCUP\_ADRD\_Falls\_Cohorts.csv} (append mode —
        safe to run one year at a time without losing prior results)
  \item \textbf{CSV is already built} — 2,495,798 rows, all 7 years (2016–2022)
\end{itemize}

\subsubsection{Table 1 Generator (\texttt{hcup\_table1.py})}
\begin{itemize}
  \item Reads from the pre-built CSV — does NOT touch raw ASC files
  \item Outputs Table 1: demographics, race, LOS, comorbidities by cohort
  \item Saves to \texttt{HCUP\_Table1.csv} and \texttt{HCUP\_Cohort\_By\_Year.csv}
\end{itemize}

\subsubsection{Report (\texttt{HCUP\_ADRD\_Report.tex})}
\begin{itemize}
  \item Answers professor's questions: cohort definitions, ADRD ICD codes,
        patient identity limitation, inclusions/exclusions, Table 1
\end{itemize}

\subsection{What Still Needs to Be Done}

\begin{tcolorbox}[colback=warnred!8, colframe=warnred,
  title=\textbf{Pending Tasks — Project 2}]
\begin{enumerate}[itemsep=6pt]
  \item \textbf{Fix PAY1 (primary payer) column positions} \\
        Current values are mostly \texttt{-9} (missing) — column byte positions in
        \texttt{get\_col\_dict()} are misaligned for PAY1. Verify against
        \texttt{SASload\_NIS\_YYYY\_Core.SAS} for each year. Once fixed, delete CSV
        and rerun \texttt{hcup\_adrd\_falls\_cohorts.py}.

  \item \textbf{Fix ZIPINC\_QRTL (zip income quartile) column positions} \\
        Same issue as PAY1 — values appear offset. Verify from SAS load file.

  \item \textbf{Fix DISPUNIFORM (discharge disposition) column positions} \\
        Values appear as 3-digit strings (e.g., \texttt{136}) instead of 1–2 digit
        codes. Column window is reading adjacent bytes. Verify and fix.

  \item \textbf{Add HOSPID to saved columns, then merge NIS Hospital file} \\
        \texttt{HOSPID} (hospital identifier) is not currently saved in the CSV.
        Steps:
        \begin{enumerate}
          \item Add \texttt{HOSPID} to the \texttt{keep} list in
                \texttt{hcup\_adrd\_falls\_cohorts.py}
          \item Delete \texttt{HCUP\_ADRD\_Falls\_Cohorts.csv} and rerun the script
                (one-time, $\sim$30–60 min)
          \item In \texttt{hcup\_table1.py}, load each year's
                \texttt{NIS\_YYYY\_Hospital.ASC} and merge on \texttt{HOSPID + YEAR}
                to get \texttt{HOSP\_REGION} and \texttt{HOSP\_LOCTEACH}
        \end{enumerate}

  \item \textbf{Resolve UNCLASSIFIED cohort (144,643 records)} \\
        Professor must decide: merge into Cohort 1, create a 4th cohort, or exclude.

  \item \textbf{Apply NIS discharge weights (\texttt{DISCWT})} \\
        For nationally representative estimates, all proportions should be weighted.
        \texttt{DISCWT} is in the NIS Core file and must be saved and applied.

  \item \textbf{Statistical testing for Table 1} \\
        Add chi-square p-values for categorical variables and ANOVA/Kruskal-Wallis
        for continuous variables comparing the three cohorts.
\end{enumerate}
\end{tcolorbox}

\subsection{Key Files — Project 2}

\begin{center}
\begin{tabular}{lp{9cm}}
\toprule
\textbf{File} & \textbf{Purpose} \\
\midrule
\texttt{hcup\_adrd\_falls\_cohorts.py} & Main cohort builder — reads raw ASC files,
  assigns cohorts, saves CSV \\
\texttt{hcup\_table1.py}              & Analysis script — reads CSV, generates Table 1 \\
\texttt{hcup\_analysis.py}            & Exploratory script — fall/Alzheimer counts (preliminary) \\
\texttt{HCUP\_ADRD\_Falls\_Cohorts.csv} & Pre-built cohort CSV — 2,495,798 rows,
  all 7 years (do not delete unless fixing column issues) \\
\texttt{HCUP\_Table1.csv}             & Table 1 output \\
\texttt{HCUP\_Cohort\_By\_Year.csv}   & Yearly cohort breakdown \\
\texttt{HCUP\_ADRD\_Report.tex}       & Report answering professor's questions \\
\bottomrule
\end{tabular}
\end{center}

\subsection{Raw Data File Locations (Do Not Move)}

\begin{center}
\begin{tabular}{ll}
\toprule
\textbf{Year} & \textbf{Path} \\
\midrule
2016 & \texttt{D:\textbackslash HCUP\textbackslash NIS\textbackslash NIS Decoded\textbackslash 2016\textbackslash NIS\_2016\_Core.ASC} \\
2017 & \texttt{D:\textbackslash HCUP\textbackslash NIS\textbackslash NIS Decoded\textbackslash 2017\textbackslash NIS\_2017\_Core.ASC} \\
2018 & \texttt{D:\textbackslash HCUP\textbackslash NIS\textbackslash NIS Decoded\textbackslash 2018\textbackslash NIS\_2018\_Core.ASC} \\
2019 & \texttt{D:\textbackslash HCUP\textbackslash NIS\textbackslash NIS Decoded\textbackslash 2019\textbackslash NIS\_2019\_Core.ASC} \\
2020 & \texttt{D:\textbackslash HCUP\textbackslash NIS\textbackslash NIS 2020\textbackslash NIS\_2020\_Core.ASC} \\
2021 & \texttt{D:\textbackslash HCUP\textbackslash NIS\textbackslash NIS 2021\textbackslash NIS\_2021\textbackslash NIS\_2021\_Core\textbackslash NIS\_2021\_Core.ASC} \\
2022 & \texttt{D:\textbackslash HCUP\textbackslash NIS\textbackslash NIS 2022\textbackslash NIS\_2022\textbackslash NIS\_2022\_Core\textbackslash NIS\_2022\_Core.ASC} \\
\bottomrule
\end{tabular}
\end{center}

\subsection{How to Run — Project 2}

\begin{enumerate}
  \item \textbf{Analysis only (CSV already built):}
        \texttt{python hcup\_table1.py} — generates Table 1 from existing CSV
  \item \textbf{Full rebuild} (if columns fixed or new years added):
        Delete \texttt{HCUP\_ADRD\_Falls\_Cohorts.csv}, then
        \texttt{python hcup\_adrd\_falls\_cohorts.py} ($\sim$30–60 min)
  \item \textbf{Do NOT run both scripts simultaneously} — the cohort builder
        appends to the CSV; running table1.py mid-build will give partial results
\end{enumerate}

\newpage

% ══════════════════════════════════════════════════════════════════════════════
\section{Summary — What Is Done vs.\ Pending}
% ══════════════════════════════════════════════════════════════════════════════

\begin{center}
\renewcommand{\arraystretch}{1.4}
\begin{tabularx}{\textwidth}{Xll}
\toprule
\textbf{Task} & \textbf{Project} & \textbf{Status} \\
\midrule
SEER data filtered and preprocessed          & SEER ML  & \textcolor{projgreen}{\textbf{Done}} \\
XGBoost model trained and calibrated         & SEER ML  & \textcolor{projgreen}{\textbf{Done}} \\
Prediction pipeline fixed (encoding bug)     & SEER ML  & \textcolor{projgreen}{\textbf{Done}} \\
Streamlit app built and tested (10 cases)    & SEER ML  & \textcolor{projgreen}{\textbf{Done}} \\
3-page LaTeX report written                  & SEER ML  & \textcolor{projgreen}{\textbf{Done}} \\
Add screenshots (out1.png–out10.png) to Overleaf & SEER ML & \textcolor{warnred}{\textbf{Pending}} \\
Guthrie external validation                  & SEER ML  & \textcolor{warnred}{\textbf{Pending — awaiting data}} \\
\midrule
HCUP raw ASC files parsed (all 7 years)      & HCUP NIS & \textcolor{projgreen}{\textbf{Done}} \\
Three cohorts built and saved to CSV         & HCUP NIS & \textcolor{projgreen}{\textbf{Done}} \\
Elixhauser comorbidity flags computed        & HCUP NIS & \textcolor{projgreen}{\textbf{Done}} \\
Table 1 (age, sex, race, comorbidities)      & HCUP NIS & \textcolor{projgreen}{\textbf{Done}} \\
Study design report for professor            & HCUP NIS & \textcolor{projgreen}{\textbf{Done}} \\
PAY1 / ZipInc / DISPUNIFORM column fix      & HCUP NIS & \textcolor{warnred}{\textbf{Pending}} \\
Hospital file merge (region, teaching)       & HCUP NIS & \textcolor{warnred}{\textbf{Pending}} \\
DISCWT survey weights applied               & HCUP NIS & \textcolor{warnred}{\textbf{Pending}} \\
Unclassified cohort decision (144K records)  & HCUP NIS & \textcolor{warnred}{\textbf{Awaiting professor}} \\
Statistical testing (chi-sq, ANOVA) for Table 1 & HCUP NIS & \textcolor{warnred}{\textbf{Pending}} \\
\bottomrule
\end{tabularx}
\end{center}

% ── Footer ───────────────────────────────────────────────────────────────────
\vspace{1cm}
\noindent\rule{\linewidth}{0.5pt}
\small
\noindent \textbf{Contact:} Sriirohit Prakash · \texttt{sprakash1@binghamton.edu} ·
Binghamton University, Department of Systems Science \& Industrial Engineering\\
\textbf{Supervisor:} Professor Ertem\\
\textbf{GitHub:} Private repository — request access from the author\\
\textbf{All raw data files} are stored on the lab machine under \texttt{D:\textbackslash HCUP\textbackslash}
(NIS) and must remain there. SEER \texttt{export.txt} is gitignored and stored locally only.

\end{document}
